%-*-tex-*-
% Copyright Michael J. Ferguson, INRS-Telecommunications
% All rights reserved. 

% ========== Equation Tags ===========
% Since equation numbering plays a big part in scientific documentation
% a special set of tag generation forms were produced for this case. 


% Special versions of autoeqnum to insert automatically insert
% parentheses around the equation number.
%  \aneq            no argument, no tag generated. This version
%                   does nothing if autonumbering is off.
%  \aneqtag{<tag>}  tag generated. This version inserts the argument
%                   inside parentheses if autoreferencing is off.

% These special forms have a \eqnumfont ... as there is no other way 
% to get the information inside. These are defined in the documentfonts.

%\let\eqnumfont=\tenrm ... for tenpoint family 

\def\aneq{\a@utotag{}{\eqnum}{\eqtagrefformat
        }{\e@qno \eqnumfont (\eqtagrefformat)}{}}
\def\aneqtag#1{\a@utotag{#1}{\eqnum}{\eqtagrefformat
        }{\e@qno\eqnumfont (\eqtagrefformat)}{\e@qno\eqnumfont (#1)}}

%===== for full left or right equation numbering and to allow ====
% \aneq and \aneqtag to work

\let\e@qql = \eqalignno
\let\l@eqql= \leqalignno
\def\e@qalignno{\let\e@qno=\relax \e@qql}
\def\l@eqalignno{\let\e@qno=\relax  \l@eqql}
%redefines \eqalignno and \leqalignno  --- normal unset mode
\let\eqalignno = \e@qalignno
\let\leqalignno = \l@eqalignno

\let\e@qno = \eqno
%---- left/right equation numbering ----
\def\leftequationnumbering{\let\eqalignno = \l@eqalignno
                         \let\leqalignno = \l@eqalignno
                         \let\e@qno = \leqno}
\def\rightequationnumbering{\let\leqalignno = \e@qalignno
                          \let\eqalignno = \e@qalignno
                         \let\e@qno = \eqno}